% sections/zh/6-conclusion.tex -- 总结、逐条局限、展望；与英文稿逐段对应。
\section{结论与展望}\label{sec:conclusion}
本文把多重度柔性作业车间调度问题的流体差距精确分解为取整损失与排序残差，并围绕这一恒等式构建进化搜索：维度由类型结构决定的配额染色体、输出保界且只面向残差的排序器，以及精确解码下的保底预选。\pending{在固定设计集、公开基准与外推档上，方法在同墙钟内达到不劣于各族最强基线的 gap，精确解码次数更少，取整差距的符号与理论预测一致。}

本研究的局限逐条对应设计边界。算例族固定 $R = 8$ 种类型、$m = 10$ 台机器且 $\bar N \le 50$，类型数的扩展未评测；换模、运输、故障、动态到达与批量拆分被排除（假设 A4 与 A5）；目标只有 makespan；数据为参数表分布与 Brandimarte 基准而非工厂数据，Hurink 集未纳入；未与长时限商用求解器及大语言模型自动设计的算法比较。\pending{官方代码不可得的学习型基线按复现协议重新实现，协议与结果一并说明。}

后续有三个方向。其一，同样的分解适用于任何在聚合意义上紧的松弛，可在批量与分批流变体上检验；其二，共享商图编码器可让排序器跨类型结构训练；其三，保底并集可接入其它候选源，如精确的小邻域搜索，而不改变接受准则。
